%=====================================================================
%  系统开发工具基础 - 第1周实验报告
%  作者：王帅
%  学号：25040032035
%  日期：2026年8月26日
%=====================================================================
\documentclass[12pt,a4paper]{article}

% ---------- 中文支持 ----------
\usepackage[UTF8]{ctex}

% ---------- 页面与排版 ----------
\usepackage[top=2.5cm, bottom=2.5cm, left=2.8cm, right=2.8cm]{geometry}
\usepackage{setspace}
\setstretch{1.4}

% ---------- 颜色与字体 ----------
\usepackage{xcolor}
\definecolor{titleblue}{RGB}{0, 51, 102}
\definecolor{sectionblue}{RGB}{0, 70, 130}
\definecolor{partpurple}{RGB}{80, 40, 120}
\definecolor{codebg}{RGB}{245, 245, 245}
\definecolor{codeframe}{RGB}{200, 200, 200}
\definecolor{codekeyword}{RGB}{0, 0, 180}
\definecolor{codecomment}{RGB}{0, 128, 0}
\definecolor{codeprompt}{RGB}{160, 32, 32}
\definecolor{linkblue}{RGB}{0, 0, 200}
\definecolor{tipgreen}{RGB}{0, 100, 0}
\definecolor{tipbg}{RGB}{240, 255, 240}

% ---------- 超链接 ----------
\usepackage{hyperref}
\hypersetup{
    colorlinks=true,
    linkcolor=titleblue,
    urlcolor=linkblue,
    citecolor=titleblue
}

% ---------- 标题页 ----------
\usepackage{titling}
\pretitle{%
    \begin{center}
    \LARGE\bfseries\color{titleblue}
    \vspace{-1cm}
    \rule{\textwidth}{1.5pt}\\[0.6cm]
}
\posttitle{%
    \\[0.4cm]
    \rule{\textwidth}{1.5pt}
    \end{center}
}
\preauthor{\begin{center}\large\color{sectionblue}}
\postauthor{\end{center}}
\predate{\begin{center}\normalsize\color{gray}}
\postdate{\end{center}}

% ---------- 代码高亮 ----------
\usepackage{listings}
\lstdefinestyle{shellstyle}{
    language=bash,
    basicstyle=\ttfamily\small,
    keywordstyle=\color{codekeyword}\bfseries,
    commentstyle=\color{codecomment}\itshape,
    stringstyle=\color{red!70!black},
    backgroundcolor=\color{codebg},
    frame=single,
    rulecolor=\color{codeframe},
    breaklines=true,
    breakatwhitespace=true,
    tabsize=4,
    showstringspaces=false,
    moredelim=[il][\color{codeprompt}]{\$},
    moredelim=[il][\color{codeprompt}]{>},
    escapeinside={(*@}{@*)},
}
\lstdefinestyle{latexstyle}{
    language=[LaTeX]TeX,
    basicstyle=\ttfamily\small,
    keywordstyle=\color{codekeyword}\bfseries,
    commentstyle=\color{codecomment}\itshape,
    stringstyle=\color{red!70!black},
    backgroundcolor=\color{codebg},
    frame=single,
    rulecolor=\color{codeframe},
    breaklines=true,
    tabsize=4,
    showstringspaces=false,
}
\lstdefinestyle{gitstyle}{
    language=bash,
    basicstyle=\ttfamily\small,
    keywordstyle=\color{codekeyword}\bfseries,
    commentstyle=\color{codecomment}\itshape,
    stringstyle=\color{red!70!black},
    backgroundcolor=\color{codebg},
    frame=single,
    rulecolor=\color{codeframe},
    breaklines=true,
    breakatwhitespace=true,
    tabsize=4,
    showstringspaces=false,
    moredelim=[il][\color{codeprompt}]{\$},
    escapeinside={(*@}{@*)},
}

% ---------- 提示框 ----------
\usepackage{tcolorbox}
\tcbuselibrary{skins, breakable}
\newtcolorbox{tipbox}[1][]{
    colback=tipbg,
    colframe=tipgreen,
    fonttitle=\bfseries,
    title={Tips},
    coltitle=tipgreen,
    breakable
}
\newtcolorbox{notebox}[1][]{
    colback=blue!3,
    colframe=sectionblue,
    fonttitle=\bfseries,
    title={Note},
    coltitle=sectionblue,
    breakable
}
\newtcolorbox{resultbox}[1][]{
    colback=yellow!3,
    colframe=orange!70,
    fonttitle=\bfseries,
    title={Result},
    coltitle=orange!70,
    breakable
}
\newtcolorbox{goalbox}[1][]{
    colback=purple!3,
    colframe=partpurple,
    fonttitle=\bfseries,
    title={Goal},
    coltitle=partpurple,
    breakable
}

% ---------- 其他 ----------
\usepackage{graphicx}
\usepackage{booktabs}
\usepackage{tabularx}
\usepackage{enumitem}
\usepackage{fancyhdr}
\pagestyle{fancy}
\fancyhf{}
\fancyhead[L]{\small\color{gray} 系统开发工具基础 · 第1周实验报告}
\fancyhead[R]{\small\color{gray} 王帅 · 25040032035}
\fancyfoot[C]{\thepage}
\renewcommand{\headrulewidth}{0.4pt}

% ---------- Part 标题格式 ----------
\usepackage{titlesec}
\titleformat{\section}
    {\normalfont\Large\bfseries\color{sectionblue}}
    {\thesection}{0.8em}
    {}[\titlerule]
\titleformat{name=\section,numberless}
    {\normalfont\Large\bfseries\color{partpurple}}
    {}{0em}
    {\centering}
    [\titlerule]

%=====================================================================
\begin{document}

% ======================== 封面 ========================
\title{系统开发工具基础 \\[0.3cm] \large 第1周实验报告}
\author{
    \textbf{王\quad 帅} \\[0.2cm]
    \normalsize 学号：25040032035 \\
    \normalsize 软件工程专业 \\
    \normalsize 中国海洋大学 计算机科学与技术学院 \\[0.2cm]
    \normalsize 授课教师：周小伟
}
\date{2026 年 8 月 26 日}

\maketitle
\vspace{-1cm}

% ======================== 目录 ========================
\tableofcontents
\newpage

% ======================== 正文 ========================

%-------------------------------------------------------
\section{实验概述}
%-------------------------------------------------------

\subsection{实验目的}

本周实验围绕三个核心主题展开，旨在建立系统开发的基础工具链认知：

\begin{enumerate}[label=\textbf{(\arabic*)}]
    \item \textbf{Shell 基础与管道}：掌握 Linux 命令行基本操作，理解文件系统导航、文件查找与权限管理，学会使用管道组合和 awk 文本处理完成复杂的日志分析任务。
    \item \textbf{Version Control \& Git}：理解版本控制的核心思想，掌握 Git 的基本工作流（add / commit / branch / merge / log），学会处理分支合并冲突。
    \item \textbf{LaTeX 文档编辑}：了解 LaTeX 排版系统的基本原理，能够编写包含中文、数学公式和表格的学术文档，并掌握交叉引用技术。
\end{enumerate}

\subsection{实验安排}

本周实验分为两个阶段：

\begin{table}[h]
\centering
\begin{tabularx}{\textwidth}{l l X}
\toprule
\textbf{阶段} & \textbf{日期} & \textbf{内容概述} \\
\midrule
\textbf{课堂实验} & 2026-08-24 & Shell 基础（mkdir + 权限管理）、Shell 管道（awk 日志分析）、Git 合并冲突、LaTeX 文档编辑 \\
\textbf{课后练习} & 2026-08-25\textasciitilde 26 & 10 个实例：Shell 基础（4）、Git 版本控制（3）、LaTeX 文档编辑（3） \\
\bottomrule
\end{tabularx}
\end{table}

\subsection{实验环境}

\begin{table}[h]
\centering
\begin{tabularx}{\textwidth}{>{\bfseries}l X}
\toprule
\textbf{项目} & \textbf{详情} \\
\midrule
操作系统 & Windows 11 \\
Linux 环境 & MSYS2 (MINGW64) \\
Git 版本 & git 2.55.0 \\
Git 客户端 & GitHub Desktop 3.6.4 \\
LaTeX 编辑器 & Overleaf（在线编译，XeLaTeX 引擎） \\
代码仓库 & \href{https://github.com/wangshuai0606/hw}{github.com/wangshuai0606/hw} \\
\bottomrule
\end{tabularx}
\end{table}

\subsection{课程资源}

\begin{itemize}[leftmargin=2em]
    \item 课程参考资料：\href{https://missing-semester-cn.github.io/}{The Missing Semester of Your CS Education（中文版）}
    \item LaTeX 参考：\href{https://oi-wiki.org/tools/latex/}{OI Wiki LaTeX 工具}
\end{itemize}

\newpage
%%%%%%%%%%%%%%%%%%%%%%%%%%%%%%%%%%%%%%%%%%%%%%%%%%%%%%%%%%%%%%%%%%%%%%%%
%%                    第一部分：课堂实验（8月24日）
%%%%%%%%%%%%%%%%%%%%%%%%%%%%%%%%%%%%%%%%%%%%%%%%%%%%%%%%%%%%%%%%%%%%%%%%
\section*{第一部分：课堂实验（2026-08-24）}
\addcontentsline{toc}{section}{第一部分：课堂实验（2026-08-24）}

本章包含 8 月 24 日课堂实验内容，共 4 个实验。这些实验在课上完成，并于当堂向助教演示讲解。

%=======================================================
\subsection{实验一：Shell 基础（mkdir + 权限管理）}
%=======================================================

\begin{goalbox}
掌握 Linux 基本文件操作命令：目录创建（\texttt{mkdir}）、文件创建（\texttt{echo}/\texttt{touch}）、文件查找（\texttt{find}）、文件复制（\texttt{cp}）以及权限管理（\texttt{chmod}）。
\end{goalbox}

\subsubsection{步骤 1：创建多层目录结构}

\begin{lstlisting}[style=shellstyle, caption={Create project directory structure}]
# Create multi-level nested directories
mkdir -p q01/input/docs
mkdir -p q01/input/tmp
mkdir -p q01/work/25040032035
\end{lstlisting}

\begin{resultbox}
成功创建以下目录树：

\texttt{q01/} → \texttt{input/} → \texttt{docs/}, \texttt{tmp/}；\texttt{work/} → \texttt{25040032035/}
\end{resultbox}

\subsubsection{步骤 2：创建各类文件}

\begin{lstlisting}[style=shellstyle, caption={Create various files}]
# Create text files
echo "Hello World" > "input/docs/notes one.txt"

# Create hidden file (starts with .)
echo "Secret!" > input/docs/.secret.txt

# Create empty file
touch input/tmp/empty.txt

# Create multi-line content file
echo -e "Started\nFinished" > input/tmp/run.log
\end{lstlisting}

\begin{tipbox}
\texttt{echo} 默认追加换行符；加 \texttt{-e} 可启用转义（如 \texttt{\textbackslash n} 换行、\texttt{\textbackslash t} 制表符）。文件名含空格时必须用引号包裹，否则 Shell 会将其视为多个参数。
\end{tipbox}

\subsubsection{步骤 3：用 find 查找并复制文件}

\begin{lstlisting}[style=shellstyle, caption={Find and copy files}]
# Find all .txt files and copy with directory structure preserved
find input -name "*.txt" | xargs -I \{\} cp --parents \{\} work/25040032035/
\end{lstlisting}

\begin{notebox}
\texttt{find} 输出文件路径列表，通过管道传给 \texttt{xargs}；\texttt{-I \{\}} 表示将每条路径替换到 \texttt{cp} 命令中；\texttt{cp -\/-parents} 会自动保留源文件的目录层级结构。
\end{notebox}

\begin{resultbox}
\texttt{work/25040032035/} 下保留了完整的目录结构，包含 \texttt{notes one.txt} 和 \texttt{.secret.txt} 两个文件。
\end{resultbox}

\subsubsection{步骤 4：设置文件权限}

\begin{lstlisting}[style=shellstyle, caption={Set file permissions}]
# Directory: owner rwx, group r-x, others ---
chmod 750 work/25040032035/input/docs

# File: owner rw-, group r--, others ---
chmod 640 work/25040032035/input/tmp/empty.txt
\end{lstlisting}

\begin{resultbox}
使用 \texttt{ls -ld} 和 \texttt{ls -l} 验证权限设置正确：\texttt{drwxr-x---} (目录) 和 \texttt{-rw-r-----} (文件)。
\end{resultbox}

\begin{tipbox}
权限数字含义：\texttt{r=4, w=2, x=1}。三个数字分别对应所有者、所属组、其他用户。例如 \texttt{750 = 7(rwx) + 5(r-x) + 0(---)}。目录的 \texttt{x} 权限表示"可进入"，没有 \texttt{x} 权限则无法 \texttt{cd} 进入该目录。
\end{tipbox}

\subsubsection{步骤 5：生成 inventory 清单}

\begin{lstlisting}[style=shellstyle, caption={Generate file inventory}]
# List all regular files with size, output to inventory
find work/25040032035 -type f -printf "\%p \%s\textbackslash n" > inventory.txt
\end{lstlisting}

\begin{resultbox}
\texttt{inventory.txt} 成功生成，每行包含文件路径和大小。
\end{resultbox}

\subsubsection{学习心得}

\begin{itemize}[leftmargin=2em]
    \item \texttt{mkdir -p} 是创建目录结构的神器，幂等操作，重复执行不会报错。
    \item \texttt{find + xargs + cp -\/-parents} 的组合非常实用，可以批量按目录结构复制文件。
    \item 权限管理中，数字模式比符号模式更简洁直观。
    \item \texttt{-printf} 是 GNU find 特有的格式化输出选项，比 \texttt{ls -l} 更适合脚本化处理。
\end{itemize}

\newpage
%=======================================================
\subsection{实验二：Shell 管道（awk 日志分析）}
%=======================================================

\begin{goalbox}
掌握管道组合和 \texttt{awk} 文本处理：通过编写 Shell 脚本分析 CSV 格式的访问日志，实现错误请求统计和平均延迟计算。
\end{goalbox}

\subsubsection{步骤 1：创建测试数据}

\begin{lstlisting}[style=shellstyle, caption={Create access.csv test data}]
cat > access.csv << 'EOF'
timestamp,method,path,status,latency
2026-08-24T08:01:12,GET,/index.html,200,45
2026-08-24T08:02:33,POST,/api/login,200,120
2026-08-24T08:03:45,GET,/api/users,500,3200
2026-08-24T08:04:01,GET,/index.html,200,38
2026-08-24T08:05:22,DELETE,/api/users/5,500,4100
2026-08-24T08:06:10,PUT,/api/users/3,200,95
2026-08-24T08:07:55,GET,/api/orders,502,5500
2026-08-24T08:08:30,GET,/index.html,200,42
EOF
\end{lstlisting}

\begin{resultbox}
成功创建包含 8 行数据的 CSV 文件，字段为：\texttt{timestamp, method, path, status, latency}。其中 3 条请求状态码 $\geq 500$（500、500、502）。
\end{resultbox}

\subsubsection{步骤 2：编写分析脚本 analyze.sh}

\begin{lstlisting}[style=shellstyle, caption={analyze.sh - Log analysis script}]
#!/bin/bash
# analyze.sh - Analyze access logs

FILE="$1"

# Check arguments and file existence
if [ -z "$FILE" ]; then
    echo "Usage: bash analyze.sh <csv-file>"
    exit 1
fi

if [ ! -f "$FILE" ]; then
    echo "Error: File '$FILE' not found"
    exit 1
fi

echo "=============================="
echo "  Log Analysis Report"
echo "  File: $FILE"
echo "=============================="

# --- 1. Count server errors (status >= 500) ---
echo ""
echo "[1] Server Errors (status >= 500) Top 5:"
echo "--------------------------------------"
awk -F',' 'NR > 1 && $4 >= 500 {
    count[$2 " " $3 " -> " $4]++
}
END {
    for (key in count)
        printf "  %-30s %d times\n", key, count[key]
}' "$FILE" | sort -t' ' -k1 -rn | head -5

# --- 2. Calculate average latency ---
echo ""
echo "[2] Average Latency:"
echo "--------------------------------------"
awk -F',' 'NR > 1 {
    sum += $5; count++
}
END {
    if (count > 0)
        printf "  Total requests: %d\n  Avg latency: %.1f ms\n  Total time:   %d ms\n",
               count, sum/count, sum
}' "$FILE"

echo ""
echo "=============================="
echo "  Analysis Complete"
echo "=============================="
\end{lstlisting}

\begin{notebox}
脚本要点：
\begin{itemize}[leftmargin=1em]
    \item \texttt{-F','}：设置 awk 的字段分隔符为逗号（CSV 格式）。
    \item \texttt{NR > 1}：跳过标题行。
    \item \texttt{\$4 >= 500}：筛选状态码大于等于 500 的服务端错误。
    \item \texttt{count[key]++}：用关联数组按请求类型分组计数。
    \item \texttt{END \{...\}}：在所有行处理完后输出汇总结果。
\end{itemize}
\end{notebox}

\subsubsection{步骤 3：运行脚本}

\begin{lstlisting}[style=shellstyle, caption={Run log analysis}]
$ bash analyze.sh access.csv
\end{lstlisting}

\begin{resultbox}
\begin{verbatim}
==============================
  Log Analysis Report
  File: access.csv
==============================

[1] Server Errors (status >= 500) Top 5:
--------------------------------------
  DELETE /api/users/5 -> 500       1 times
  GET /api/orders -> 502           1 times
  GET /api/users -> 500            1 times

[2] Average Latency:
--------------------------------------
  Total requests: 8
  Avg latency: 1655.0 ms
  Total time:   13240 ms

==============================
  Analysis Complete
==============================
\end{verbatim}
\end{resultbox}

\subsubsection{学习心得}

\begin{itemize}[leftmargin=2em]
    \item \texttt{awk} 是 Unix 文本处理的瑞士军刀，尤其擅长按列处理结构化数据。
    \item 脚本编写中参数检查是好习惯，可以让脚本更健壮。
    \item 管道 + \texttt{sort} + \texttt{head} 的组合可以轻松实现"Top N"排行。
\end{itemize}

\newpage
%=======================================================
\subsection{实验三：Git 合并冲突}
%=======================================================

\begin{goalbox}
掌握 Git 分支管理和冲突解决：通过手动制造并解决合并冲突，深入理解 Git 的分支模型和三路合并机制。
\end{goalbox}

\subsubsection{步骤 1：初始化仓库并配置}

\begin{lstlisting}[style=gitstyle, caption={Initialize Git repository}]
# Create project directory and init
mkdir ~/hw/q_merge && cd ~/hw/q_merge
git init

# Configure user info
git config user.name "wangshuai"
git config user.email "wangshuai@example.com"
\end{lstlisting}

\subsubsection{步骤 2：创建 main 分支并提交初始内容}

\begin{lstlisting}[style=gitstyle, caption={Create initial commit}]
# Create hello.py
cat > hello.py << 'EOF'
def greet(name):
    return f"Hello, {name}! Welcome to Git."

if __name__ == "__main__":
    print(greet("World"))
EOF

git add hello.py
git commit -m "Initial commit: add hello.py"
\end{lstlisting}

\subsubsection{步骤 3：创建 feature 分支并修改}

\begin{lstlisting}[style=gitstyle, caption={Modify on feature branch}]
# Create and switch to feature branch
git checkout -b feature

# Modify greet function with decoration
cat > hello.py << 'EOF'
def greet(name):
    msg = f"Hello, {name}! Welcome to Git."
    return "=" * 40 + "\n" + msg + "\n" + "=" * 40

if __name__ == "__main__":
    print(greet("World"))
EOF

git add hello.py
git commit -m "feature: add decoration border"
\end{lstlisting}

\subsubsection{步骤 4：切回 main 修改同一位置 → 产生冲突}

\begin{lstlisting}[style=gitstyle, caption={Modify same location on main}]
# Switch back to main
git checkout main

# Modify the same function differently
cat > hello.py << 'EOF'
def greet(name, lang="en"):
    if lang == "zh":
        return f"Ni hao, {name}!"
    return f"Hello, {name}! Welcome to Git."

if __name__ == "__main__":
    print(greet("World", lang="zh"))
    print(greet("World"))
EOF

git add hello.py
git commit -m "main: add multi-language support"
\end{lstlisting}

\subsubsection{步骤 5：合并并解决冲突}

\begin{lstlisting}[style=gitstyle, caption={Merge and resolve conflict}]
# Try to merge feature branch -> CONFLICT!
git merge feature
# Output: CONFLICT (content): Merge conflict in hello.py

# View conflict markers
cat hello.py
# Shows: <<<<<<< HEAD ... ======= ... >>>>>>> feature

# Manually resolve: combine both features
cat > hello.py << 'EOF'
def greet(name, lang="en"):
    if lang == "zh":
        msg = f"Ni hao, {name}!"
    else:
        msg = f"Hello, {name}! Welcome to Git."
    return "=" * 40 + "\n" + msg + "\n" + "=" * 40

if __name__ == "__main__":
    print(greet("World", lang="zh"))
    print(greet("World"))
EOF

# Mark conflict resolved and commit
git add hello.py
git commit -m "Merge feature: combine decoration and multi-language"
\end{lstlisting}

\begin{resultbox}
冲突成功解决！最终代码同时包含了两个分支的功能：多语言支持 + 装饰边框效果。
\end{resultbox}

\begin{lstlisting}[style=gitstyle, caption={View merge history}]
$ git log --oneline --graph

*   abc1234 (HEAD -> main) Merge feature branch
|\
| * def5678 (feature) feature: add decoration border
* | ghi9012 main: add multi-language support
|/
* jkl3456 Initial commit: add hello.py
\end{lstlisting}

\subsubsection{学习心得}

\begin{tipbox}
\textbf{冲突解决三步法：}
\begin{enumerate}[leftmargin=1em]
    \item \textbf{看懂标记}：\texttt{<<<<<<< HEAD} 到 \texttt{=======} 是当前分支内容，\texttt{=======} 到 \texttt{>>>>>>> feature} 是待合并分支内容。
    \item \textbf{手动编辑}：删除冲突标记，保留/融合需要的代码。
    \item \textbf{提交解决}：\texttt{git add} + \texttt{git commit}，Git 会自动生成合并提交记录。
\end{enumerate}
\end{tipbox}

\newpage
%=======================================================
\subsection{实验四：LaTeX 文档编辑（公式与表格）}
%=======================================================

\begin{goalbox}
掌握 LaTeX 数学公式排版（\texttt{equation} 环境）和表格排版（\texttt{table}/\texttt{tabular} 环境），并学会使用交叉引用技术。
\end{goalbox}

\subsubsection{步骤 1：编写求和公式（equation 环境）}

\begin{lstlisting}[style=latexstyle, caption={Summation formula example}]
\documentclass[12pt,a4paper]{article}
\usepackage{amsmath}

\begin{document}

% Gauss summation formula
\begin{equation}
    \label{eq:gauss-sum}
    S = \sum_{i=1}^{n} i = \frac{n(n+1)}{2}
\end{equation}

As shown in equation~(\ref{eq:gauss-sum}),
the sum of first $n$ positive integers equals $\frac{n(n+1)}{2}$.

\end{document}
\end{lstlisting}

\begin{resultbox}
编译成功！公式~(\ref{eq:gauss-sum})自动编号，交叉引用~\texttt{\textbackslash ref} 正确显示公式编号。
\end{resultbox}

\subsubsection{步骤 2：编写工具对比表格}

\begin{lstlisting}[style=latexstyle, caption={Tool comparison table}]
\documentclass[12pt,a4paper]{article}
\usepackage{booktabs}

\begin{document}

\begin{table}[h]
\centering
\caption{Comparison of Text Editing Tools}
\label{tab:tool-compare}
\begin{tabular}{lccc}
\toprule
\textbf{Feature} & \textbf{LaTeX} & \textbf{Word} & \textbf{Markdown} \\
\midrule
Formula typesetting  & 5 stars & 3 stars & 2 stars \\
Table typesetting    & 4 stars & 4 stars & 2 stars \\
Learning curve       & Steep   & Easy    & Easy    \\
Version control      & 5 stars & 1 star  & 4 stars \\
Collaboration        & 2 stars & 4 stars & 3 stars \\
\bottomrule
\end{tabular}
\end{table}

As shown in Table~\ref{tab:tool-compare}, LaTeX excels in formula typesetting
and version control friendliness.

\end{document}
\end{lstlisting}

\begin{resultbox}
编译成功！三线表（\texttt{booktabs} 风格）排版精美，交叉引用~\texttt{\textbackslash ref\{tab:tool-compare\}} 正确显示编号。
\end{resultbox}

\subsubsection{步骤 3：交叉引用综合示例}

\begin{lstlisting}[style=latexstyle, caption={Cross-reference comprehensive example}]
\documentclass[12pt,a4paper]{article}
\usepackage{amsmath}
\usepackage{booktabs}

\begin{document}

\section{Experiment Content}

\subsection{Math Formula}
\begin{equation}
    \label{eq:sum}
    S = \sum_{i=1}^{n} i = \frac{n(n+1)}{2}
\end{equation}

\subsection{Tool Comparison}
\begin{table}[h]
\centering
\caption{Text Editing Tools}
\label{tab:tools}
\begin{tabular}{lcc}
\toprule
\textbf{Feature} & \textbf{LaTeX} & \textbf{Word} \\
\midrule
Formula typesetting & 5 stars & 3 stars \\
Version control     & 5 stars & 1 star  \\
\bottomrule
\end{tabular}
\end{table}

\section{Summary}
As shown in Eq.~(\ref{eq:sum}) and Table~\ref{tab:tools},
LaTeX is indispensable for academic typesetting.

\end{document}
\end{lstlisting}

\begin{notebox}
交叉引用的关键点：
\begin{itemize}[leftmargin=1em]
    \item \texttt{\textbackslash label\{key\}}：在公式/表格内部定义标签。
    \item \texttt{\textbackslash ref\{key\}}：在正文中引用标签，自动生成编号。
    \item 公式引用常用 \texttt{\textasciitilde(\textbackslash ref\{...\})} 格式（带括号）。
    \item 表格引用常用 \texttt{表\textasciitilde\textbackslash ref\{...\}} 格式。
    \item \texttt{\textasciitilde}（不可断空格）确保"表"字和编号不会被分行。
\end{itemize}
\end{notebox}

\subsubsection{学习心得}

\begin{itemize}[leftmargin=2em]
    \item \texttt{equation} 环境会自动编号，配合 \texttt{\textbackslash label} + \texttt{\textbackslash ref} 实现交叉引用，修改后编号自动更新。
    \item \texttt{booktabs} 的三线表比默认边框表格美观得多，是学术论文的标准做法。
    \item LaTeX 的"标记式"写作方式初学有些别扭，但一旦熟练后效率极高。
\end{itemize}

\newpage
%%%%%%%%%%%%%%%%%%%%%%%%%%%%%%%%%%%%%%%%%%%%%%%%%%%%%%%%%%%%%%%%%%%%%%%%
%%                    第二部分：课后练习（10个实例）
%%%%%%%%%%%%%%%%%%%%%%%%%%%%%%%%%%%%%%%%%%%%%%%%%%%%%%%%%%%%%%%%%%%%%%%%
\section*{第二部分：课后练习 --- 10 个实例}
\addcontentsline{toc}{section}{第二部分：课后练习 --- 10 个实例}

本次课后练习共完成 10 个实例，覆盖 Shell 基础（4 个）、Git 版本控制（3 个）和 LaTeX 文档编辑（3 个）。所有代码已上传至 GitHub 仓库，每个实例单独 commit。

%=======================================================
\subsection{Shell 基础（实例 1--4）}
%=======================================================

%-------------------------------------------------------
\subsubsection{实例 1：mkdir --- 创建目录与文件}

\textbf{目标：}学会使用 \texttt{mkdir} 和 \texttt{touch} 创建多层目录结构与文件。

\begin{lstlisting}[style=shellstyle, caption={Create project directory structure}]
# Create homework directories
mkdir -p ~/hw/q01/input/docs
cd ~/hw/q01

# Create files in input/docs
touch input/docs/readme.txt
touch input/docs/notes.txt
touch input/docs/data.csv
\end{lstlisting}

\begin{resultbox}
目录 \texttt{q01/input/docs/} 下成功创建了 3 个文件，目录结构完整。
\end{resultbox}

\begin{tipbox}
\texttt{mkdir -p} 可以一次性创建多级目录，如果父目录已存在不会报错，非常实用。
\end{tipbox}

%-------------------------------------------------------
\subsubsection{实例 2：find --- 文件查找}

\textbf{目标：}掌握 \texttt{find} 命令按名称、类型、大小等条件搜索文件。

\begin{lstlisting}[style=shellstyle, caption={Search files with find}]
# Find all .txt files
find ~/hw/q01 -name "*.txt"

# Find all regular files
find ~/hw/q01 -type f

# Find files larger than 0 bytes
find ~/hw/q01 -type f -size +0c
\end{lstlisting}

\begin{resultbox}
成功找到 \texttt{readme.txt} 和 \texttt{notes.txt} 两个文件，以及 \texttt{data.csv}。
\end{resultbox}

\begin{notebox}
\texttt{find} 的 \texttt{-name} 参数支持通配符（如 \texttt{*.txt}），而 \texttt{-type f} 限定只搜索普通文件，\texttt{-type d} 则搜索目录。
\end{notebox}

%-------------------------------------------------------
\subsubsection{实例 3：grep --- 文本搜索}

\textbf{目标：}学会使用 \texttt{grep} 在文件中搜索匹配的文本行。

\begin{lstlisting}[style=shellstyle, caption={Search text with grep}]
# Prepare data files
echo -e "apple\nbanana\ncherry\napple pie" > ~/hw/q03/fruits.txt
echo -e "cat\ndog\nbird\ndog fish" > ~/hw/q03/animals.txt

# Search for "apple"
grep "apple" ~/hw/q03/fruits.txt

# Search for "dog"
grep "dog" ~/hw/q03/animals.txt

# Recursive search
grep -r "apple" ~/hw/q03/
\end{lstlisting}

\begin{resultbox}
\texttt{grep "apple"} 匹配到 2 行（apple、apple pie）；\texttt{grep "dog"} 匹配到 2 行（dog、dog fish）。
\end{resultbox}

%-------------------------------------------------------
\subsubsection{实例 4：管道与 sort/uniq}

\textbf{目标：}理解管道（\texttt{|}）的工作原理，学会将多个命令组合使用。

\begin{lstlisting}[style=shellstyle, caption={Pipeline commands}]
# Find txt files, sort and count
find ~/hw/q03 -name "*.txt" | sort | uniq -c

# Merge files, sort and deduplicate
cat ~/hw/q03/fruits.txt ~/hw/q03/animals.txt | sort | uniq
\end{lstlisting}

\begin{resultbox}
管道成功将 \texttt{find} 的输出传递给 \texttt{sort} 排序，再由 \texttt{uniq} 去重统计。每条内容各出现 1 次。
\end{resultbox}

\begin{tipbox}
管道是 Unix 哲学的精髓——每个工具只做一件事，通过管道串联就能完成复杂任务。就像乐高积木，简单的零件组合出强大的功能！
\end{tipbox}

\newpage
%=======================================================
\subsection{Git 版本控制（实例 5--7）}
%=======================================================

%-------------------------------------------------------
\subsubsection{实例 5：git branch --- 创建与切换分支}

\textbf{目标：}理解分支的概念，学会创建和切换分支。

\begin{lstlisting}[style=shellstyle, caption={Branch operations}]
# Create q05 sub-project
cd ~/hw
mkdir q05 && cd q05
git init

# Create initial file and commit
echo "line1" > file.txt
git add .
git commit -m "Initial commit"

# Create and switch to feature branch
git branch feature
git checkout feature

# Modify file on feature branch
echo "line2-feature" >> file.txt
git add .
git commit -m "feature: add line2"
\end{lstlisting}

\begin{resultbox}
成功创建 \texttt{feature} 分支并切换到该分支，提交了独立修改。
\end{resultbox}

\begin{notebox}
\texttt{git branch <name>} 只创建分支不切换；\texttt{git checkout <name>} 切换到指定分支。也可以用 \texttt{git checkout -b <name>} 一步完成创建+切换。
\end{notebox}

%-------------------------------------------------------
\subsubsection{实例 6：git merge --- 合并与冲突解决}

\textbf{目标：}学会合并分支并手动解决合并冲突。

\begin{lstlisting}[style=shellstyle, caption={Merge and resolve conflicts}]
# Switch to master and modify same location
git checkout master
echo "line2-master" >> file.txt
git add .
git commit -m "master: add line2"

# Try merge feature -> CONFLICT!
git merge feature
# Output: CONFLICT (content): Merge conflict in file.txt

# View conflict
cat file.txt
# Shows: <<<<<<< HEAD / ======= / >>>>>>> feature

# Resolve: keep both line2 entries
# Final file.txt content:
# line1
# line2-master
# line2-feature

git add .
git commit -m "Resolved: keep both line2"
\end{lstlisting}

\begin{resultbox}
冲突成功解决！最终文件保留了两个分支各自添加的 \texttt{line2} 内容。
\end{resultbox}

\begin{tipbox}
冲突标记的含义：\texttt{<<<<<<< HEAD} 到 \texttt{=======} 是当前分支的内容，\texttt{=======} 到 \texttt{>>>>>>> feature} 是待合并分支的内容。解决冲突就是手动决定保留哪些内容。
\end{tipbox}

%-------------------------------------------------------
\subsubsection{实例 7：git log --- 查看提交历史}

\textbf{目标：}使用 \texttt{git log} 查看提交历史，理解项目的提交脉络。

\begin{lstlisting}[style=shellstyle, caption={View commit history}]
# Full commit history
git log

# Concise one-line mode
git log --oneline

# Graphical branch view
git log --oneline --graph
\end{lstlisting}

\begin{resultbox}
输出结果清晰展示了分支与合并的历史：
\begin{verbatim}
*   abc1234 (HEAD -> master) Resolved: keep both line2
|\
| * def5678 (feature) feature: add line2
* | ghi9012 master: add line2
|/
* jkl3456 Initial commit
\end{verbatim}
\end{resultbox}

\begin{notebox}
\texttt{--graph} 参数用 ASCII 字符画出分支树，非常直观。\texttt{--oneline} 只显示前 7 位 commit hash 和标题，简洁高效。
\end{notebox}

\newpage
%=======================================================
\subsection{LaTeX 文档编辑（实例 8--10）}
%=======================================================

以下三个实例均在 \textbf{Overleaf} 在线编辑器中完成，使用 \texttt{XeLaTeX} 引擎编译以支持中文。每个实例的源码已单独提交至 GitHub 仓库。

%-------------------------------------------------------
\subsubsection{实例 8：基本文档结构}

\textbf{目标：}掌握 LaTeX 文档的基本结构，包括导言区、正文区、标题、章节等。

\begin{lstlisting}[style=latexstyle, caption={Basic document structure}]
\documentclass[12pt,a4paper]{article}
\usepackage[UTF8]{ctex}

\title{My First LaTeX Document}
\author{Shuai Wang}
\date{August 26, 2026}

\begin{document}
\maketitle

\section{Introduction}
This is a basic LaTeX document example.

\section{Document Structure}
LaTeX documents consist of preamble and body:
\begin{itemize}
    \item Preamble: document class, packages
    \item Body: all sections and content
\end{itemize}

\subsection{Common Commands}
\begin{enumerate}
    \item \textbackslash section\{\} -- Section heading
    \item \textbackslash textbf\{\} -- Bold text
    \item \textbackslash textit\{\} -- Italic text
\end{enumerate}

\end{document}
\end{lstlisting}

\begin{resultbox}
编译成功！生成了包含标题页、章节编号、有序/无序列表的完整 PDF 文档。
\end{resultbox}

%-------------------------------------------------------
\subsubsection{实例 9：数学公式}

\textbf{目标：}学会在 LaTeX 中输入行内公式、独立公式、多行对齐公式、矩阵等。

\begin{lstlisting}[style=latexstyle, caption={Math formula examples}]
\documentclass[12pt,a4paper]{article}
\usepackage[UTF8]{ctex}
\usepackage{amsmath}

\begin{document}

% Inline formula
Mass-energy equivalence: $E = mc^2$

% Display formula
$$\int_{-\infty}^{\infty} e^{-x^2} dx = \sqrt{\pi}$$

% Aligned equations
\begin{align}
    (a+b)^2 &= a^2 + 2ab + b^2 \\
    (a-b)^2 &= a^2 - 2ab + b^2
\end{align}

% Matrix
$$A = \begin{pmatrix}
    1 & 2 & 3 \\
    4 & 5 & 6 \\
    7 & 8 & 9
\end{pmatrix}$$

\end{document}
\end{lstlisting}

\begin{resultbox}
编译成功！数学公式排版精美，包括行内公式、高斯积分、完全平方公式组、$3\times3$ 矩阵等均正确渲染。
\end{resultbox}

\begin{tipbox}
\texttt{amsmath} 宏包是 LaTeX 数学排版的利器，提供 \texttt{align}、\texttt{gather}、\texttt{cases} 等强大环境。写数学论文必备！
\end{tipbox}

%-------------------------------------------------------
\subsubsection{实例 10：表格}

\textbf{目标：}学会使用 \texttt{tabular} 环境和 \texttt{booktabs} 宏包制作专业表格。

\begin{lstlisting}[style=latexstyle, caption={Table examples}]
\documentclass[12pt,a4paper]{article}
\usepackage[UTF8]{ctex}
\usepackage{booktabs}

\begin{document}

% Basic table with borders
\begin{table}[h]
\centering
\begin{tabular}{|l|c|r|}
\hline
Name & Age & Score \\
\hline
Alice & 18 & 92 \\
Bob   & 19 & 85 \\
\hline
\end{tabular}
\caption{Student Score Table}
\end{table}

% Three-line table (academic style)
\begin{table}[h]
\centering
\begin{tabular}{lcr}
\toprule
Course & Credits & Score \\
\midrule
Calculus & 5 & 92 \\
English  & 4 & 85 \\
\bottomrule
\end{tabular}
\caption{Course Score Table}
\end{table}

\end{document}
\end{lstlisting}

\begin{resultbox}
编译成功！生成了两种风格的表格：带边框的基础表格和学术论文常用的三线表（\texttt{booktabs} 风格）。
\end{resultbox}

\begin{notebox}
\texttt{booktabs} 宏包提供的 \texttt{\textbackslash toprule}、\texttt{\textbackslash midrule}、\texttt{\textbackslash bottomrule} 比普通的 \texttt{\textbackslash hline} 更美观，是学术论文的推荐做法。
\end{notebox}

\newpage
%-------------------------------------------------------
\section{GitHub 仓库与提交记录}
%-------------------------------------------------------

\subsection{仓库信息}

\begin{table}[h]
\centering
\begin{tabularx}{\textwidth}{>{\bfseries}l X}
\toprule
\textbf{项目} & \textbf{内容} \\
\midrule
GitHub 用户名 & wangshuai0606 \\
仓库名称 & hw \\
仓库地址 & \href{https://github.com/wangshuai0606/hw}{https://github.com/wangshuai0606/hw} \\
可见性 & Public \\
\bottomrule
\end{tabularx}
\end{table}

\subsection{Commit 记录汇总}

\begin{table}[h]
\centering
\begin{tabularx}{\textwidth}{cl l X}
\toprule
\textbf{序号} & \textbf{Commit} & \textbf{时间} & \textbf{说明} \\
\midrule
1 & 3cb0316 & 2026-08-26 & 实例1-4: Shell 基础 (mkdir, find, grep, pipe) \\
2 & c2aa75a & 2026-08-26 & 实例5-6: Git 分支与合并冲突练习 \\
3 & b43564a & 2026-08-26 & 实例8: LaTeX 基本文档结构 \\
4 & 64f6f14 & 2026-08-26 & 实例9: 数学公式 \\
5 & dd12680 & 2026-08-26 & 实例10: 表格 \\
\bottomrule
\end{tabularx}
\end{table}

\begin{notebox}
所有代码均为单独 commit 提交，符合老师"禁止单次全量提交"的要求。完整的 commit 记录截图见 GitHub 仓库的 Commits 页面。
\end{notebox}

\subsection{仓库截图}

\begin{figure}[h]
\centering
\includegraphics[width=0.9\textwidth]{screenshot_repo_home}
\caption{GitHub 仓库首页}
\end{figure}

\begin{figure}[h]
\centering
\includegraphics[width=0.9\textwidth]{screenshot_commits}
\caption{Commits 提交记录}
\end{figure}

\newpage
%-------------------------------------------------------
\section{实验总结与心得}
%-------------------------------------------------------

\subsection{知识点总结}

\begin{table}[h]
\centering
\begin{tabularx}{\textwidth}{l l X}
\toprule
\textbf{主题} & \textbf{核心命令/概念} & \textbf{收获} \\
\midrule
Shell 基础 & mkdir, find, chmod, cp & 文件操作与权限管理是系统开发的基石 \\
Shell 管道 & awk, sort, head, pipe & 管道组合 + awk 处理结构化数据威力巨大 \\
Git & branch, merge, conflict & 分支管理让协作开发井然有序，冲突不可怕 \\
LaTeX & equation, tabular, label/ref & 专业排版的学术利器，公式和表格特别强大 \\
\bottomrule
\end{tabularx}
\end{table}

\subsection{遇到的问题与解决方案}

\begin{enumerate}[leftmargin=2em]
    \item \textbf{问题：} 在 Windows CMD 中运行 Linux 命令（如 \texttt{rm -rf}、\texttt{cd \textasciitilde/hw}）报错。\\
    \textbf{解决：} Windows CMD 不支持 Linux 语法，需使用 MSYS2 (MINGW64) 终端，或使用对应的 Windows 命令。

    \item \textbf{问题：} Overleaf 使用 pdfLaTeX 编译中文文档报错（CTeX fontset `fandol' is unavailable）。\\
    \textbf{解决：} 将编译器切换为 XeLaTeX，完美支持中文。

    \item \textbf{问题：} Git 合并分支时出现冲突，不知道如何处理。\\
    \textbf{解决：} 通过 \texttt{cat file.txt} 查看冲突标记，手动编辑保留所需内容，然后 \texttt{git add} + \texttt{git commit} 完成解决。

    \item \textbf{问题：} \texttt{awk} 脚本中处理 CSV 文件时，标题行也被当作数据处理。\\
    \textbf{解决：} 使用 \texttt{NR > 1} 条件跳过第一行（标题行），只处理数据行。

    \item \textbf{问题：} LaTeX 交叉引用编译后显示 "??" 而非正确编号。\\
    \textbf{解决：} 需要编译两次——第一次生成 \texttt{.aux} 文件记录标签，第二次才能正确引用。
\end{enumerate}

\subsection{心得体会}

通过本周的实验，我对系统开发工具链有了全新的认识：

\begin{itemize}[leftmargin=2em]
    \item \textbf{Shell 是程序员的"手"}：以前总觉得图形界面更方便，但学会命令行后发现，批量操作、自动化脚本都离不开 Shell。\texttt{find + xargs} 的组合可以完成复杂的文件批量处理，\texttt{awk} 更是文本分析的神器。管道的设计思想——每个工具做一件事并组合使用——非常优雅。

    \item \textbf{Git 是程序的"记忆"}：版本控制让每次修改都有迹可循，分支机制让多人协作互不干扰。冲突虽然看起来可怕，但理解三路合并的原理后，解决冲突就变成了简单的文本编辑工作。

    \item \textbf{LaTeX 是程序的"脸面"}：作为理工科学生，以后写论文、报告都需要排版。LaTeX 的学习曲线虽然陡峭，但一旦上手，排版效果远超 Word，尤其是数学公式的渲染堪称完美。交叉引用的自动编号功能更是让长文档维护变得轻松。
\end{itemize}

\subsection{后续学习计划}

\begin{enumerate}[leftmargin=2em]
    \item 深入学习 Shell 脚本编程，包括变量、条件判断（\texttt{if/else}）、循环（\texttt{for/while}）、函数等。
    \item 学习 Git 远程协作：\texttt{push}、\texttt{pull}、\texttt{fetch}、Pull Request 工作流。
    \item 在 Overleaf 上尝试更复杂的 LaTeX 排版：bibliography 参考文献管理、图片插入、多文件项目组织。
\end{enumerate}

\vspace{1cm}
\begin{center}
\color{gray}\rule{0.6\textwidth}{0.5pt}\\[0.3cm]
{\small --- 第 1 周实验报告 · 完 ---}
\end{center}

\end{document}
